% AI-written, AI cross-reviewed; no human sign-off; not a claim that RH is proved
% RH Notes blueprint source excerpts for Z1.5.06
% Verbatim excerpts in the notation of the parent blueprint. The labelled editions and source ranges appear below.
% See the unit YAML for the current proof status. Earlier-round passages retain their original wording.
\section*{Z1.5.06. Rational Frobenius trace bound}

% Parent source SHA-256: 1af64922b98ed2310744411b1f18b25cc426c97220646321471836e1a3a409d6
% blueprint/ch_P.tex lines 12399 through 12418; r5:Z1.5.06
\subsection{Z1.5.06. A rational eigenvalue classifier gives the rank-three traces}\label{r5:Z1.5.06}
From $\mathrm{CountQAlt}$ and $\mathrm{NewAdmAlt}$ obtain constants $c,A$ independent of the number field $K$, positive $e,M$, prime $p\nmid M$, place $w\mid p$, and $h\ge0$. For these data one finite index set of size $\le c[M(e+1)]^A$ classifies all $\mathrm{Adm2B}$ representations at inverse Frobenius of residue power $p^h$ for each odd $\ell\ne p$, allowed closure, $\iota$ and place over $w$. A classifier must separate nonisomorphic representations within each fixed admissible family. Its complex values are permuted by field automorphisms, bounded by $\sqrt{p^h}+2$, and become algebraic integers after multiplication by $p^{he}$.

{\footnotesize \nolinkurl{UpperTheory.GalN.V.Win.N.frobenius_trace_count_bound_of_rat_eigensystem_count_alt} (\nolinkurl{Development/AutomorphicArithmetic/AutomorphicVarieties/ModularityArgument/ThresholdArguments/DenseThresholdCases/FrobeniusTraceCountBoundOfCmEigensystemCount.lean}\allowbreak:348--398). \nolinkurl{UpperTheory.GalN.V.Win.EigData}, \nolinkurl{TrData}, \nolinkurl{TrData.Classifies} (\nolinkurl{Development/AutomorphicArithmetic/Foundations/GaloisRepresentations/AdmissibleResidues/WinDefs.lean}\allowbreak:18--38,70--90).}

Apply the rank-two instance of $\mathrm{CountQAlt}$ at the rational place $v$ below $w$, weight $(e,-e)$ and purity parameter zero. The weight is strictly decreasing because $e\ge1$, its reversed sum vanishes, and both absolute values are at most $e$. On the resulting finite set, replace the multiset of two eigenvalues $\{b_1,b_2\}$ by
\[
 t=1+b_1^{-h}+b_2^{-h}.
\]
Retain the original labels, including when two values of $t$ coincide. This construction commutes with automorphisms of $\mathbb C$. Let $R=\sqrt{p^h}$ and $u_i=|b_i^{-h}|$. The weight-zero window and determinant norm give $0<u_i\le R$ and $u_1u_2=1$. Thus $(R-u_1)(R-u_2)\ge0$ yields $u_1+u_2\le R+1/R\le R+1$, proving the stated bound on $t$. Each $p^e b_i^{-1}$ is integral by the inverse-eigenvalue bound. Its $h$th power and the integer $p^{he}$ prove integrality for $p^{he}t$. All formulas include $h=0$.

{\footnotesize \nolinkurl{UpperTheory.GalN.V.Win.traces_from_eigenvalues_rank2} (\nolinkurl{Development/AutomorphicArithmetic/Foundations/GaloisRepresentations/AdmissibleResidues/FrobeniusTraceCountBoundOfRatEigensystemCount.lean}\allowbreak:152--230). \nolinkurl{place_over_rational_prime}, \nolinkurl{frobenius_trace_count_bound_of_rat_eigensystem_count_alt} (\nolinkurl{Development/AutomorphicArithmetic/AutomorphicVarieties/ModularityArgument/ThresholdArguments/DenseThresholdCases/FrobeniusTraceCountBoundOfCmEigensystemCount.lean}\allowbreak:342--367).}

For the chosen place $B$ above $w$, the contracted prime in $\mathbb Z$ is $(p)$ and restriction gives a place above $v$. Lift the residue-field $p$-Frobenius to $\tau$ in its decomposition group: on the residue of the integral closure, finite-field Frobenius is an automorphism and the residue-action surjectivity used in Section~\ref{r6:review:Z1.5.06} lifts that automorphism to the stabilizer of the prime. The lifted automorphism fixes the valuation ring, and its congruence on integral elements extends to fractions whose denominator is a unit there. On the residue field $\tau^h$ and $\sigma|_{\mathbb Q}$ both act by $x\mapsto x^{p^h}$, so their quotient lies in inertia. For a rank-three $r$ the $\mathrm{Adm2B}$ data give a rational rank-two $W$ with $r$ conjugate to its restriction on the first and third coordinates and identity on the middle one. $\mathrm{NewAdmAlt}$ puts $W$ in the rational classifier. Unramifiedness kills this inertia quotient, hence $W((\sigma^{-1})|_{\mathbb Q})=W(\tau^{-1})^h$. The characteristic polynomial of $W(\tau^{-1})$ has roots $b_1^{-1},b_2^{-1}$ after $\iota$. Cayley--Hamilton gives the two-term recurrence for the traces of its powers, with initial traces $2,b_1^{-1}+b_2^{-1}$, yielding the displayed $t$ upon adjoining the middle line. Equal labels give a rational intertwiner for $W,W'$; its block extension, conjugated back by the two chosen bases, intertwines the original representations. The cardinality bound follows because the index set was retained.

{\footnotesize \nolinkurl{ReciprocityTheory.Gal.exists_frob_over} (\nolinkurl{ReciprocityTheory/GaloisValuations/GlobalB.lean}\allowbreak:238--310). \nolinkurl{UpperTheory.frobenius_pow_resDeg_place_norm} (\nolinkurl{Development/AutomorphicArithmetic/AutomorphicVarieties/BuildingsResidueChars/FrobeniusPowResDegPlaceNorm.lean}\allowbreak:14--25,43--93). \nolinkurl{UpperTheory.map_trace_eq_sum_roots_of_charpoly_map} (\nolinkurl{Development/AutomorphicArithmetic/AutomorphicVarieties/BuildingsResidueChars/MapTraceEqSumRootsOfCharpolyMap.lean}\allowbreak:5--49). \nolinkurl{UpperTheory.block_conjugate_rank_three_rep} (\nolinkurl{Development/AutomorphicArithmetic/AutomorphicVarieties/BuildingsResidueChars/BlockConjugateRankThreeRep.lean}\allowbreak:5--43). \nolinkurl{UpperTheory.GalN.V.Win.N.frobenius_trace_count_bound_of_rat_eigensystem_count_alt} (\nolinkurl{Development/AutomorphicArithmetic/AutomorphicVarieties/ModularityArgument/ThresholdArguments/DenseThresholdCases/FrobeniusTraceCountBoundOfCmEigensystemCount.lean}\allowbreak:368--398).}

\emph{Complete for this transfer.} [Hypothesis, discharged at the application:] the rational counting input is $\mathrm{countQ\_alt}$ applied to the automorphic, level and restriction data. Newform admissibility is the explicit $\mathrm{NewAdmAlt}$ hypothesis of its consumer, supplied there by $\mathrm{newAdmS\_alt}$ from the published inputs.

{\footnotesize \nolinkurl{UpperTheory.GalN.V.Win.N.numberU3U_of_adm2A_imp_adm2B}, \nolinkurl{numberU3U_of_prints} (\nolinkurl{Development/AutomorphicArithmetic/AutomorphicVarieties/ModularityArgument/ThresholdArguments/DenseThresholdCases/KeepNumU.lean}\allowbreak:25--32,79--80).}

% blueprint/ch_P.tex lines 12582 through 12585; r6:review:Z1.5.06
\subsection{Z1.5.06. Frobenius modulo inertia}\label{r6:review:Z1.5.06} \emph{Original argument on p.~\pageref{r5:Z1.5.06}.} 

The rational classifier supplies a label for each admissible $W$. For every prime of the closure above the chosen rational place it asserts both that inertia acts trivially and that the roots at \emph{inverse} Frobenius are the inverses of the entries of the labelled multiset. In the integral closure the norm-$p$ Frobenius on the residue field lifts to the stabilizer of the prime by surjectivity of the residue-action map for a profinite group. Here the residue extension is algebraic over a finite field, so the $p$-power endomorphism is an automorphism; the surjectivity being invoked is precisely the profinite statement used at line 261 below. The lift extends the residue congruence from integral elements to their local fractions. Thus its $h$-th power and the restriction of $\sigma$ differ by inertia at this same place. In inverse Frobenius notation the roots contributing to the rank-three trace are $b_1^{-h},b_2^{-h}$ and $1$. The argument uses the stated rational eigenvalue classifier and admissibility comparison.
\par{\footnotesize \nolinkurl{UpperTheory.GalN.V.Win.EigData.Classifies} (\nolinkurl{Development/AutomorphicArithmetic/Foundations/GaloisRepresentations/AdmissibleResidues/WinDefs.lean}\allowbreak:34--38). \nolinkurl{UpperTheory.V.Af.FrobEig} (\nolinkurl{Development/AutomorphicArithmetic/AutomorphicVarieties/AutomorphicForms/BlochKatoSelmerVanishes.lean}\allowbreak:69--72). \nolinkurl{ReciprocityTheory.Gal.exists_frob_over} (\nolinkurl{ReciprocityTheory/GaloisValuations/GlobalB.lean}\allowbreak:238--310). \nolinkurl{UpperTheory.GalN.V.Win.N.frobenius_trace_count_bound_of_rat_eigensystem_count_alt} (\nolinkurl{Development/AutomorphicArithmetic/AutomorphicVarieties/ModularityArgument/ThresholdArguments/DenseThresholdCases/FrobeniusTraceCountBoundOfCmEigensystemCount.lean}\allowbreak:348--398).\par}

